\documentclass[12pt]{article}
\usepackage[T1]{fontenc}
\usepackage[utf8]{inputenc}
\usepackage{lmodern}
\usepackage{textcomp}
\usepackage{enumitem}
\usepackage{geometry}
\geometry{margin=1in}
\begin{document}
\begin{center}
Answer Key - Philosophy - Undergraduate Level Assessment 6 \
\small (Answers are ordered exactly as in the question paper.)
\end{center}

\section*{Multiple Choice}
\begin{enumerate}[label=\arabic*.]
\item \textbf{Answer:} C
\\ \textit{Options:} A) elements of nature that do not exist independently; B) only things existing apart from our minds; C) only sensations existing in our minds; D) physical objects; E) manifestations of our subconscious
\\ \textit{Note:} Berkeley argues that heat and cold are not intrinsic qualities of objects but rather subjective sensations experienced in the mind, highlighting his idealist philosophy that emphasizes the role of perception in understanding reality.
\item \textbf{Answer:} A
\\ \textit{Options:} A) to implement stricter guidelines for classroom discussions.; B) to promote greater understanding of historical and contemporary oppression.; C) to increase the number of safe spaces on campus.; D) to mandate sensitivity training for all students.; E) to increase funding for mental health services on campus.
\\ \textit{Note:} Lukianoff and Haidt argue that implementing stricter guidelines for classroom discussions encourages open dialogue and critical thinking, countering the tendency towards vindictive protectiveness that stifles diverse perspectives.
\item \textbf{Answer:} C
\\ \textit{Options:} A) students to confabulate reasons.; B) the pursuit of justice by marking out racism, sexism, and classism.; C) labeling, by assigning global negative traits to persons.; D) universities to bear overly burdensome legal obligations.
\\ \textit{Note:} Lukianoff and Haidt argue that the trend of identifying microaggressions fosters a culture of labeling individuals with broad negative characteristics, which oversimplifies complex human behavior and undermines constructive dialogue.
\item \textbf{Answer:} B
\\ \textit{Options:} A) no valid conclusion can be drawn; B) the conclusion must affirm the consequent; C) the conclusion must deny the consequent; D) the conclusion must deny the antecedent
\\ \textit{Note:} In a hypothetical syllogism, if the minor premise affirms the antecedent, it logically leads to the conclusion affirming the consequent due to the structure of conditional reasoning, which states that if "If P, then Q" is true and P is affirmed, Q must also be affirmed.
\item \textbf{Answer:} A
\\ \textit{Options:} A) its rarity or frequency.; B) societal norms and values.; C) the individual's personal preference.; D) the amount of effort required to obtain it.; E) the potential pain that might accompany it.
\\ \textit{Note:} According to Mill, the value of a particular pleasure is influenced by its rarity or frequency because he believes that the quality and intensity of pleasure can vary based on how often it is experienced, with rarer pleasures often being more valuable due to their limited availability.
\end{enumerate}

\section*{True / False}
\begin{enumerate}[label=\arabic*.]
\setcounter{enumi}{5}
\item \textbf{Answer:} True
\\ \textit{Options:} A) True; B) False
\\ \textit{Note:} Butler argues that all other rights stem from the fundamental right to justice, which he considers the basis of moral and ethical considerations in human interactions.
\item \textbf{Answer:} True
\\ \textit{Options:} A) True; B) False
\\ \textit{Note:} Singer's principle of equality asserts that all beings capable of suffering deserve equal consideration of their interests, thereby guiding the ethical treatment of humans by recognizing their equal moral worth.
\item \textbf{Answer:} False
\\ \textit{Options:} A) True; B) False
\\ \textit{Note:} Ashford's article argues that meeting everyone's basic needs is not only possible but also essential for social justice and ethical considerations, challenging the notion of impossibility.
\item \textbf{Answer:} False
\\ \textit{Options:} A) True; B) False
\\ \textit{Note:} The term Guru-Panth refers to the collective community of Sikhs following the teachings of the Gurus, rather than an individual's personal spiritual journey.
\item \textbf{Answer:} False
\\ \textit{Options:} A) True; B) False
\\ \textit{Note:} Hobbes argues that good is defined by individual desire rather than universal consensus, as each person has their own subjective preferences influenced by their circumstances and experiences.
\end{enumerate}

\section*{Long Form Response}
\begin{enumerate}[label=\arabic*.]
\setcounter{enumi}{10}
\item \textbf{Answer:} Hursthouse critiques the assertion that an adequate moral theory must offer clear, actionable guidance for individuals, deeming the claim controversial. She argues that moral theories should not be reduced to mere decision-making tools that provide straightforward answers to ethical dilemmas. Instead, Hursthouse emphasizes the importance of moral understanding and the development of virtuous character, suggesting that moral reasoning often involves nuanced considerations that cannot be distilled into simple directives. This perspective challenges the notion that moral theories must prioritize clear guidance, highlighting the complexity of moral life and the role of individual judgment. Ultimately, Hursthouse advocates for a richer conception of morality that accommodates ambiguity and the cultivation of virtue rather than merely seeking actionable conclusions.
\\ \textit{Note:} correct because it accurately reflects Hursthouse's critique of the expectation for moral theories to deliver clear-cut guidance, instead emphasizing the role of moral understanding and character development in ethical reasoning.
\item \textbf{Answer:} Pogge argues that individuals in wealthy countries have moral responsibilities towards the global poor, which can be understood through the lens of libertarian duties. He contends that while libertarianism traditionally emphasizes individual rights and freedoms, it also necessitates a recognition of the harm caused by systemic injustices that perpetuate poverty. Wealthy individuals benefit from unjust global structures that maintain inequality, which means they have a duty to rectify these injustices. Therefore, libertarian duties extend beyond mere non-interference to include positive obligations to alleviate the suffering of the impoverished, highlighting a moral imperative for the affluent to contribute to global justice and equity. This perspective challenges the conventional libertarian view that personal wealth should remain untouched, advocating instead for a more socially responsible approach to wealth distribution.
\\ \textit{Note:} The answer correctly highlights Pogge's argument that wealthy individuals have moral responsibilities towards the global poor by linking libertarian duties to the acknowledgment of harm caused by systemic injustices and the obligation to amend the inequalities perpetuated by unjust global structures.
\end{enumerate}

\vfill
\noindent\textit{End of Answer Key}
\end{document}